% Part: first-order-logic
% Chapter: sequent-calculus
% Section: proof-theoretic-notions

\documentclass[../../../include/open-logic-section]{subfiles}

\begin{document}

\begin{editorial}
  এই অংশে সিকোয়েন্ট কলনের প্রমাণযোগ্যতা-সম্পর্ক ও সঙ্গতির সংজ্ঞাগুলি
  একত্র করা হয়েছে।
  % BN-SRC-027 TARGET-CORRECTION-BEGIN
  \par\smallskip\noindent\textnormal{\small\textbf{উৎস-সংশোধন (BN-SRC-027):} হিমায়িত ইংরেজি সম্পাদকীয় নির্দেশে এই সংজ্ঞাগুলিকে “natural deduction”-এর বলা হয়েছে। ফাইলটি সিকোয়েন্ট কলন অধ্যায়ের অংশ, পরের সংজ্ঞাগুলি প্রত্যেকেই $\Log{LK}$ সিকোয়েন্ট ও তার বিধি ব্যবহার করে, এবং স্বাভাবিক নিষ্পাদনের পৃথক সংজ্ঞা আগের ভাগে আছে। তাই বাংলা নির্দেশে “সিকোয়েন্ট কলন” পুনরুদ্ধার করা হয়েছে।} % BN-SRC-027 TARGET-CORRECTION-END
\end{editorial}

\iftag{FOL}
      {\olfileid{fol}{seq}{ptn}}
      {\olfileid{pl}{seq}{ptn}}

\olsection{প্রমাণতাত্ত্বিক ধারণা}

\begin{explain}
যেভাবে আমরা কয়েকটি গুরুত্বপূর্ণ অর্থতাত্ত্বিক ধারণা (বৈধতা,
অনুসিদ্ধান্ত, পরিতৃপ্তিযোগ্যতা) সংজ্ঞায়িত করেছি, এবার তেমনি তাদের
অনুরূপ \emph{প্রমাণতাত্ত্বিক ধারণা} সংজ্ঞায়িত করি। এগুলি কোনো
!!{structure}s-তে !!{sentence}s-এর পরিতৃপ্তির সাহায্যে সংজ্ঞায়িত নয়;
বরং কয়েকটি সিকোয়েন্টের !!{derivability} বা !!{nonderivability}-র
সাহায্যে সংজ্ঞায়িত। এই ধারণাগুলি যে পরস্পরের সঙ্গে মিলে যায়, তা এক
গুরুত্বপূর্ণ আবিষ্কার। তাদের মিলই \emph{বিশুদ্ধতা} ও \emph{পূর্ণতা
উপপাদ্য}-র বক্তব্য।
\end{explain}

\begin{defn}[উপপাদ্য]
কোনো !!^a{sentence}~$!A$ একটি \emph{উপপাদ্য}, যদি
$\quad \Sequent !A$ সিকোয়েন্টটির $\Log{LK}$-তে কোনো
!!a{derivation} থাকে। $!A$ উপপাদ্য হলে $\Proves !A$ এবং উপপাদ্য না
হলে $\Proves/ !A$ লিখি।
\end{defn}

\begin{defn}[!!^{derivability}]
কোনো !!^a{sentence}~$!A$ ঠিক তখনই !!{sentence}s-এর সেট~$\Gamma$
\emph{থেকে !!{derivable}}, অর্থাৎ $\Gamma \Proves !A$, যখন একটি
সসীম উপসেট~$\Gamma_0 \subseteq \Gamma$ এবং $\Gamma_0$-র
!!{sentence}s-এর একটি অনুক্রম~$\Gamma_0'$ আছে, যেন $\Log{LK}$
$\Gamma_0' \Sequent !A$ !!{derive}s করে। $!A$ যদি $\Gamma$ থেকে
!!{derivable} না হয়, তবে $\Gamma \Proves/ !A$ লিখি।
\end{defn}

সংকোচন, দুর্বলীকরণ ও অদলবদল-বিধির জন্য~$\Gamma_0'$-এ
!!{sentence}s-এর ক্রম ও সংখ্যা গুরুত্বপূর্ণ নয়: $\Gamma_0'
\Sequent !A$ যদি !!{derivable} হয়, তবে~$\Gamma_0'$-এর মতো একই
!!{sentence}s-যুক্ত যেকোনো~$\Gamma_0''$-এর জন্য $\Gamma_0''
\Sequent !A$-ও নিষ্পাদনযোগ্য। যেমন, $\Gamma_0 = \{!B, !C\}$ হলে
$\Gamma_0' = \tuple{!B, !B, !C}$ এবং $\Gamma_0'' =
\tuple{!C, !C, !B}$—দুটিই শুধু~$\Gamma_0$-র !!{sentence}s-যুক্ত
অনুক্রম। একটিকে নিয়ে গঠিত সিকোয়েন্টটি !!{derivable} হলে অন্যটিও তাই;
যেমন:
\begin{prooftree}
  \AxiomC{}
  \Deduce$!B, !B, !C \fCenter !A$
  \RightLabel{\LeftR{\Contraction}}
  \UnaryInf$!B, !C \fCenter !A$
  \RightLabel{\LeftR{\Exchange}}
  \UnaryInf$!C, !B \fCenter !A$
  \RightLabel{\LeftR{\Weakening}}
  \UnaryInf$!C, !C, !B \fCenter !A$
\end{prooftree}
এখন থেকে~$\Gamma_0$ যদি !!{sentence}s-এর একটি সসীম সেট হয়, তবে
$\Gamma_0 \Sequent !A$ দিয়ে এমন যেকোনো সিকোয়েন্ট বোঝাব যার পূর্বাংশ
হলো~$\Gamma_0$-র !!{sentence}s-এর একটি অনুক্রম; দরকার হলে সংকোচন,
অদলবদল ও দুর্বলীকরণ নিঃশব্দে অন্তর্ভুক্ত বলে ধরব।

\begin{defn}[সঙ্গতি]
বাক্যের কোনো সেট~$\Gamma$ ঠিক তখনই \emph{অসঙ্গত}, যখন এমন একটি সসীম
উপসেট~$\Gamma_0 \subseteq \Gamma$ আছে যেন $\Log{LK}$
$\Gamma_0 \Sequent \quad$ !!{derive}s করে। $\Gamma$ অসঙ্গত না হলে,
অর্থাৎ প্রত্যেক সসীম~$\Gamma_0 \subseteq \Gamma$-র জন্য $\Log{LK}$
$\Gamma_0 \Sequent \quad$ !!{derive} করে না হলে, সেটিকে \emph{সঙ্গত}
বলি।
\end{defn}

\begin{prop}[প্রতিবিম্ব ধর্ম]
\ollabel{prop:reflexivity}
$!A \in \Gamma$ হলে $\Gamma \Proves !A$।
\end{prop}

\begin{proof}
প্রারম্ভিক সিকোয়েন্ট $!A \Sequent !A$ !!{derivable}, এবং
$\{!A\} \subseteq \Gamma$।
\end{proof}

\begin{prop}[একঘেয়েতা]
\ollabel{prop:monotonicity}
$\Gamma \subseteq \Delta$ এবং $\Gamma \Proves !A$ হলে
$\Delta \Proves !A$।
\end{prop}

\begin{proof}
ধরা যাক $\Gamma \Proves !A$; অর্থাৎ এমন একটি সসীম $\Gamma_0
\subseteq \Gamma$ আছে, যেন $\Gamma_0 \Sequent !A$ !!{derivable}।
$\Gamma \subseteq \Delta$ বলে $\Gamma_0$-ও~$\Delta$-র একটি সসীম
উপসেট। সুতরাং $\Gamma_0 \Sequent !A$-এর !!{derivation}-টিই
$\Delta \Proves !A$ দেখায়।
\end{proof}

\begin{prop}[পরিযায়িতা]
\ollabel{prop:transitivity}
$\Gamma \Proves !A$ এবং $\{!A\} \cup \Delta \Proves !B$ হলে
$\Gamma \cup \Delta \Proves !B$।
\end{prop}

\begin{proof}
$\Gamma \Proves !A$ হলে একটি সসীম $\Gamma_0 \subseteq \Gamma$ এবং
$\Gamma_0 \Sequent !A$-এর কোনো !!a{derivation}~$\pi_0$ আছে।
$\{!A\} \cup \Delta \Proves !B$ হলে কোনো সসীম উপসেট
$\Delta_0 \subseteq \Delta$-র জন্য $!A, \Delta_0 \Sequent !B$-এর
কোনো !!a{derivation}~$\pi_1$ আছে। নিচের !!{derivation}-টি বিবেচনা করো:
\begin{prooftree}
  \AxiomC{}
  \RightLabel{$\pi_0$}
  \Deduce$\Gamma_0 \fCenter !A$
  \AxiomC{}
  \RightLabel{$\pi_1$}
  \Deduce$!A, \Delta_0 \fCenter !B$
  \RightLabel{\Cut}
  \BinaryInf$\Gamma_0, \Delta_0 \fCenter !B$
\end{prooftree}
$\Gamma_0 \cup \Delta_0 \subseteq \Gamma \cup \Delta$ বলে এটি
$\Gamma \cup \Delta \Proves !B$ দেখায়।
\end{proof}

লক্ষ করো, এর বিশেষ ক্ষেত্র হিসেবে $\Gamma \Proves !A$ এবং
$!A \Proves !B$ হলে $\Gamma \Proves !B$। আবার, $!A_1,
\dots, !A_n \Proves !B$ এবং প্রত্যেক~$i$-এর জন্য $\Gamma \Proves !A_i$
হলে $\Gamma \Proves !B$।

\begin{prop}
\ollabel{prop:incons}
$\Gamma$ ঠিক তখনই অসঙ্গত, যখন প্রত্যেক বাক্য~$!A$-এর জন্য
$\Gamma \Proves {!A}$।
\end{prop}

\begin{proof}
অনুশীলনী।
\end{proof}

\begin{prob}
\olref[fol][seq][ptn]{prop:incons} প্রমাণ করো।
\end{prob}

\begin{prop}[সংহতি]
  \ollabel{prop:proves-compact}
  \begin{enumerate}
  \item $\Gamma \Proves !A$ হলে এমন একটি সসীম উপসেট $\Gamma_0
    \subseteq \Gamma$ আছে যেন $\Gamma_0 \Proves !A$।
  \item $\Gamma$-র প্রত্যেক সসীম উপসেট সঙ্গত হলে $\Gamma$ সঙ্গত।
  \end{enumerate}
\end{prop}

\begin{proof}
  \begin{enumerate}
    \item $\Gamma \Proves !A$ হলে এমন একটি সসীম উপসেট
      $\Gamma_0 \subseteq \Gamma$ আছে, যেন $\Gamma_0
      \Sequent !A$ সিকোয়েন্টটির কোনো !!a{derivation} আছে। ফলে
      $\Gamma_0 \Proves !A$।
    \item $\Gamma$ অসঙ্গত হলে এমন একটি সসীম উপসেট
      $\Gamma_0 \subseteq \Gamma$ আছে, যেন $\Log{LK}$
      $\Gamma_0 \Sequent \quad$ !!{derive}s করে। কিন্তু তখন~$\Gamma_0$
      হলো~$\Gamma$-র একটি অসঙ্গত সসীম উপসেট।
  \end{enumerate}
\end{proof}

\end{document}
